\section{Experimental Design}

The experiments ask whether a recommendation can be made credible by routing
candidate proposals through the same evidence and safety protocol. We separate
five questions: whether invalid or leaky evidence is rejected, whether the
selection logic can recover known observation structure in controlled tasks,
whether proposer orderings reach useful candidates under fixed budgets, whether
the structured proposer protocol is stable across LLM providers, and what the
resulting framework recommends in a real FluSurv-NET case study.
Table~\ref{tab:rqs} summarizes the design and the interpretation each evidence
mode can support.

\begin{table*}[t]
\centering
\small
\caption{Research-question driven experimental design. Each row uses the same
candidate schema, verifier, no-leakage rule, and claim-boundary audit, but
supports a different level of interpretation.}
\label{tab:rqs}
\begin{tabularx}{\textwidth}{p{0.9cm}p{3.3cm}p{4.0cm}X}
\toprule
RQ & Question & Evidence & Supported interpretation \\
\midrule
RQ1 & Can the verifier block unsafe evidence? &
Negative-set verifier tests, prompt/no-leakage audits, claim audits &
Evidence provenance, leakage control, and claim-boundary reliability. \\
RQ2 & Can known observation and delay labels be recovered? &
Generic structured time-series tasks with direct, lagged, mixture, and proxy
observations &
Software validation of observation-label and delay-label selection logic. \\
RQ3 & Do proposer orderings reach useful candidates under fixed budgets? &
Synthetic candidate execution, frozen real-data replay, and bounded real-data
execution &
Candidate-budget efficiency, not new real-data forecasting superiority. \\
RQ4 & Is the proposer protocol provider-agnostic? &
OpenAI/GPT, Claude, Gemini, and DeepSeek under the same schema, allowlist, and
verifier &
Structured proposer reliability and replay efficiency, not provider ranking
for clinical forecasting. \\
RQ5 & What does the framework recommend in FluSurv-NET? &
Frozen benchmark, discovery ablations, paired rolling-origin comparisons, and
a compact multi-season appendix &
Age- and objective-dependent case-study recommendations. \\
\bottomrule
\end{tabularx}
\end{table*}

\subsection{No-leakage split and evidence modes}

All selection decisions use non-test evidence only. A candidate may be proposed
from the series name, the candidate allowlist, budget state, verifier feedback,
non-final validation or rolling-origin summaries, candidate complexity, and
numerical-risk indicators. Prompts, feedback objects, and policy inputs are
audited for forbidden fields, including held-out test MAE, test winners, test
ranks, post-selection test comparisons, and paper recommendations derived from
test evidence. Held-out test MAE is computed or read only after a model has
already been selected and is stored under explicitly post-hoc fields such as
\texttt{post\_selection\_test\_mae}. This rule is deliberately stricter than a
typical benchmark table because the central object of study is the
recommendation protocol rather than a single model score.

The paper uses four evidence modes. First, synthetic structured recovery
provides ground-truth observation and delay labels under generic numerical
time-series tasks. Second, frozen replay reorders existing compact benchmark
rows without refitting models. Third, bounded execution evaluates a small
allowlisted union of real-data candidates with reduced deterministic fitting
settings and then deletes temporary model artifacts. Fourth, the frozen
FluSurv-NET discovery-ablation package provides the main real-data case study.
These modes are intentionally not interchangeable: synthetic recovery supports
software validation, replay supports ordering efficiency, bounded execution
supports budget-limited execution feasibility, and the FluSurv-NET benchmark
supports age- and objective-dependent case-study conclusions.

\subsection{Real-data case study}

The real-data case study uses CDC FluSurv-NET/RESP-NET hospitalization-rate
series for Overall, \texttt{0-4 yr}, \texttt{5-17 yr}, \texttt{18-49 yr},
\texttt{50-64 yr}, and \texttt{>=65 yr}
\citep{cdcFluSurvNet2026,cdcRespNet2026,naquin2024flusurvnet}. The frozen
discovery-ablation package compares standard forecasting baselines, manual
mechanistic baselines, constrained observation-aware structured search, and
same-grammar ablations. We keep the frozen benchmark fixed and treat all
additional selection-policy, agentic-proposer, and bounded-execution experiments
as downstream audits or budget-efficiency evaluations. No downstream stage
modifies the frozen discovery-ablation artifacts.

The compact multi-season robustness appendix uses completed seasons
2018--19, 2019--20, 2021--22, 2022--23, 2023--24, and 2024--25 as independent
within-season trajectories. The incomplete 2020--21 season and preliminary
2025--26 season are excluded. This appendix is not a previous-season to
future-season transfer forecast; it is a compact season-separated caveat for
whether the pediatric observation-aware signal is stable across completed
seasons under a reduced budget.

\subsection{Candidate families, ablations, and policies}

The evaluated candidate space contains four main families. Forecasting
baselines include last observed, rolling-mean, and small automatic ARIMA
models. Manual mechanistic baselines include deterministic compartmental
templates with fixed structure. Structured-search candidates use the bounded
grammar described above, including observation-aware search over direct,
delayed, mixture, and proxy observation labels. Ablation candidates remove or
alter parts of the selection protocol: no-observation-search fixes the
observation label, validation-only selection ignores the stability-aware
score, no-stability selection removes the stability penalty, and
random/exhaustive search provide same-grammar controls.

The main policy is Pareto-epsilon selection over non-test evidence:
validation score when available, rolling-origin MAE, complexity, numerical
risk, and family diversity. Weighted scoring is retained as a deterministic
ablation, and the hard-veto decision tree provides a simple interpretable
policy that can prefer a baseline when it is within epsilon of a more complex
candidate. The policy layer is evaluated because model recommendation is not
only a question of the best held-out score: in a partially observed setting,
the credibility of a recommendation also depends on whether simpler baselines
are competitive, whether the observation map is supported by ablations, and
whether numerical failures or flagged rows are being used as evidence.

\subsection{Verifier and claim-boundary tests}

The verifier is tested with a negative set containing missing required fields,
duplicate candidate identifiers, absolute artifact paths, invalid observation
or delay labels, out-of-allowlist model names, held-out test metrics included
as selection evidence, and numerical-failure rows used as positive support.
The same no-leakage audit is applied to prompt contexts, feedback objects,
selection inputs, and compact outputs in the candidate-execution and
provider-benchmark experiments. Claim-boundary auditing then maps the evidence
into allowed and rejected statements. In particular, the auditor can allow
claims about proposal quality, budget efficiency, synthetic recovery, and
age/objective-dependent case-study behavior, while rejecting claims about
state-of-the-art forecasting, FluSight leaderboard performance, autonomous
science, provider clinical superiority, and real-world mechanism recovery.

\subsection{Synthetic structured recovery}

The synthetic sweep is a generic structured time-series validation, not a
biological simulation. It uses five tasks with known labels: direct signal,
one-week lagged signal, two-week lagged signal, mixture observation, and
hidden-component proxy. The observation formulas are
\begin{align}
\mathrm{direct}:&\quad y_t=z_t+\epsilon_t,\\
\mathrm{lagged}(d):&\quad y_t=z_{t-d}+\epsilon_t,\quad d\in\{1,2\},\\
\mathrm{mixture}:&\quad y_t=\alpha z_t+(1-\alpha)u_t+\epsilon_t,\\
\mathrm{proxy}:&\quad y_t=u_t+\epsilon_t.
\end{align}
The expanded sweep evaluates 100 seeds, five noise levels
$\{0,0.025,0.05,0.10,0.15\}$, five candidate budgets
$\{3,5,10,20,40\}$, and six policies or baselines: Pareto-epsilon, weighted
score, hard-veto decision tree, deterministic seed proposer, random-label
baseline, and no-observation-label baseline. The resulting compact table
contains 75,000 policy-condition rows. The primary metrics are observation
label recovery, delay label recovery, candidate-family recovery, rolling
error, budget to recover the true label, valid proposal rate, duplicate rate,
and top-epsilon hit rate.

\subsection{Candidate-budget execution and replay}

Candidate-budget experiments simulate the question a user faces under limited
evaluation budget: after $k$ candidates have become available, has the
proposer exposed a useful candidate? The synthetic execution layer actually
scores candidates on the generic toy tasks and contains 24,000 evaluated
rows. The frozen real-data replay layer uses existing compact FluSurv-NET
rows only and contains 64 replay rows. The bounded real-data execution layer
evaluates a small allowlisted candidate union on Overall, \texttt{0-4 yr},
\texttt{18-49 yr}, and \texttt{>=65 yr}. It uses reduced deterministic fitting
settings, executes each unique series--model pair once, writes temporary
outputs under a disposable root, suppresses diagnostic plot writes, summarizes
the metrics, and deletes temporary artifacts. The committed bounded execution
summary contains 60 series-budget rows, 32 unique real-data executions, and a
52-row no-leakage audit in which all rows pass.

For real-data replay and bounded execution, selection uses rolling or
validation evidence, complexity, numerical risk, and verifier status. Test MAE
is attached only after selection and is interpreted descriptively. The
comparison includes deterministic seed, random candidate ordering,
failure-guided ordering, mock agent ordering, and an oracle reference ordering
derived from the full compact candidate set. The oracle is a reference for
budget-to-top-epsilon, not a deployable proposer.

\subsection{Constrained agentic proposer and cross-provider benchmark}

The language-model component is evaluated as a constrained proposer only. Each
provider receives the same task protocol: an explicit candidate allowlist,
series metadata, budget and round index, prior verifier feedback, duplicate
candidate information, family and observation-label coverage, and non-test
rolling or validation evidence when available. Each provider must return JSON
candidate records with fields for candidate id, family, model name,
observation label, delay label, intended role, rationale, expected failure
mode, and paired-ablation target. The output is parsed, allowlist-checked,
verified, and audited before replay. Provider output cannot generate model
code, add candidate families, alter the allowlist, execute forecasts, or see
held-out test evidence.

We first ran a single-provider LLM-backed proposal smoke test to check proposal validity and
safety, then a three-round iterative proposer loop in frozen replay. The final
cross-provider benchmark compares OpenAI/GPT, Anthropic Claude, Google Gemini,
and DeepSeek as structured proposers. It uses six series, three repeats, three
rounds, three candidates per round, and budgets $\{3,6,9\}$. Both iterative
and single-shot provider modes are compared to deterministic seed,
failure-guided, random, and oracle-reference orderings. The benchmark reports
schema parse success, verifier-valid proposal rate, duplicate rate,
out-of-allowlist rejection, family and observation-label diversity,
between-repeat stability, selected-model agreement, and frozen-replay
top-epsilon hit. It is not a provider leaderboard for clinical forecasting.
